\documentclass[10pt,a4paper]{amsart}
\usepackage{amsmath,amssymb,amsthm,graphicx}
\usepackage{verbatim}
\newenvironment{reference}{\comment}{\endcomment}
%\newenvironment{reference}{}{}
\newenvironment{slogan}{\comment}{\endcomment}
\newenvironment{history}{\comment}{\endcomment}

% For commutative diagrams we use Xy-pic
\usepackage[all]{xy}

% We use 2cell for 2-commutative diagrams.
\xyoption{2cell}
% Theorem environments.
%
\theoremstyle{plain}
\newtheorem{theorem}[subsection]{Theorem}
\newtheorem{proposition}[subsection]{Proposition}
\newtheorem{lemma}[subsection]{Lemma}

\theoremstyle{definition}
\newtheorem{definition}[subsection]{Definition}
\newtheorem{example}[subsection]{Example}
\newtheorem{exercise}[subsection]{Exercise}
\newtheorem{situation}[subsection]{Situation}

\theoremstyle{remark}
\newtheorem{remark}[subsection]{Remark}
\newtheorem{remarks}[subsection]{Remarks}

\numberwithin{equation}{subsection}

% Macros
%
\def\lim{\mathop{\mathrm{lim}}\nolimits}
\def\colim{\mathop{\mathrm{colim}}\nolimits}
\def\Spec{\mathop{\mathrm{Spec}}}
\def\Hom{\mathop{\mathrm{Hom}}\nolimits}
\def\Ext{\mathop{\mathrm{Ext}}\nolimits}
\def\SheafHom{\mathop{\mathcal{H}\!\mathit{om}}\nolimits}
\def\SheafExt{\mathop{\mathcal{E}\!\mathit{xt}}\nolimits}
\def\Sch{\mathit{Sch}}
\def\Mor{\mathop{\mathrm{Mor}}\nolimits}
\def\Ob{\mathop{\mathrm{Ob}}\nolimits}
\def\Sh{\mathop{\mathit{Sh}}\nolimits}
\def\NL{\mathop{N\!L}\nolimits}
\def\CH{\mathop{\mathrm{CH}}\nolimits}
\def\proetale{{pro\text{-}\acute{e}tale}}
\def\etale{{\acute{e}tale}}
\def\QCoh{\mathit{QCoh}}
\def\Ker{\mathop{\mathrm{Ker}}}
\def\Im{\mathop{\mathrm{Im}}}
\def\Coker{\mathop{\mathrm{Coker}}}
\def\Coim{\mathop{\mathrm{Coim}}}

% Boxtimes
%
\DeclareMathSymbol{\boxtimes}{\mathbin}{AMSa}{"02}

%
% Macros for moduli stacks/spaces
%
\def\QCohstack{\mathcal{QC}\!\mathit{oh}}
\def\Cohstack{\mathcal{C}\!\mathit{oh}}
\def\Spacesstack{\mathcal{S}\!\mathit{paces}}
\def\Quotfunctor{\mathrm{Quot}}
\def\Hilbfunctor{\mathrm{Hilb}}
\def\Curvesstack{\mathcal{C}\!\mathit{urves}}
\def\Polarizedstack{\mathcal{P}\!\mathit{olarized}}
\def\Complexesstack{\mathcal{C}\!\mathit{omplexes}}
% \Pic is the operator that assigns to X its picard group, usage \Pic(X)
% \Picardstack_{X/B} denotes the Picard stack of X over B
% \Picardfunctor_{X/B} denotes the Picard functor of X over B
\def\Pic{\mathop{\mathrm{Pic}}\nolimits}
\def\Picardstack{\mathcal{P}\!\mathit{ic}}
\def\Picardfunctor{\mathrm{Pic}}
\def\Deformationcategory{\mathcal{D}\!\mathit{ef}}
\usepackage{cleveref}
\crefname{originallemma}{original lemma}{original lemmas}
\Crefname{originallemma}{Original lemma}{Original lemmas}
\makeatletter
\newlabel{more-morphisms-lemma-geometrically-connected-fibres-towards-normal}{{0AY8}{source}{}{sourcepointer-1}{}}
\newlabel{more-morphisms-lemma-geometrically-connected-fibres-towards-normal@cref}{{[originallemma][0][]0AY8}{[1][0][]source}{}{}{}}
\makeatother
\crefname{originalsection}{original section}{original sections}
\Crefname{originalsection}{Original section}{Original sections}
\makeatletter
\newlabel{morphisms-section-dimension-formula}{{02JT}{source}{}{sourcepointer-2}{}}
\newlabel{morphisms-section-dimension-formula@cref}{{[originalsection][0][]02JT}{[1][0][]source}{}{}{}}
\makeatother
\crefname{originalsection}{original section}{original sections}
\Crefname{originalsection}{Original section}{Original sections}
\makeatletter
\newlabel{morphisms-section-nagata}{{0359}{source}{}{sourcepointer-3}{}}
\newlabel{morphisms-section-nagata@cref}{{[originalsection][0][]0359}{[1][0][]source}{}{}{}}
\makeatother
\crefname{originalsection}{original section}{original sections}
\Crefname{originalsection}{Original section}{Original sections}
\makeatletter
\newlabel{morphisms-section-normalization}{{035E}{source}{}{sourcepointer-4}{}}
\newlabel{morphisms-section-normalization@cref}{{[originalsection][0][]035E}{[1][0][]source}{}{}{}}
\makeatother
\crefname{originallemma}{original lemma}{original lemmas}
\Crefname{originallemma}{Original lemma}{Original lemmas}
\makeatletter
\newlabel{varieties-lemma-dimension-locally-algebraic}{{0A21}{source}{}{sourcepointer-5}{}}
\newlabel{varieties-lemma-dimension-locally-algebraic@cref}{{[originallemma][0][]0A21}{[1][0][]source}{}{}{}}
\makeatother
\crefname{originallemma}{original lemma}{original lemmas}
\Crefname{originallemma}{Original lemma}{Original lemmas}
\makeatletter
\newlabel{varieties-lemma-modification-normal-iso-over-codimension-1}{{0BFP}{source}{}{sourcepointer-6}{}}
\newlabel{varieties-lemma-modification-normal-iso-over-codimension-1@cref}{{[originallemma][0][]0BFP}{[1][0][]source}{}{}{}}
\makeatother
\title{A proper birational morphism of a two dimensional Noetherian Nagata scheme is dominated by normalized blowups at closed points}
\author{The Stacks Project Authors}
\date{stacks-project:resolve.tex:local-snapshot-20261012; GNU FDL 1.2 or later}
\begin{document}
\maketitle
\noindent\textit{Editorial note.} All delivered mathematical text is the author's own, copied from the pinned Stacks Project chapter file resolve.tex; no statement or proof step was rewritten, reordered or removed.\par\medskip
\noindent\textit{Editorial note.} The counted proof is the author's complete argument: the reduction to normal surfaces with the finiteness of the bad points, the induction on the number of contracted curves with the injection comparing bad curves along a normalized blowup, and Zariski's argument producing a common exceptional factor that strictly decreases the orders and hence kills one bad curve after finitely many steps.\par\medskip
\noindent\textit{Editorial note.} The label lemma-dominate-by-normalized-blowing-up names the counted statement, and the proof environment immediately following it in the source is delivered verbatim as its proof.\par\medskip
\noindent\textit{Editorial note.} The counted statement is carried by the author's own lemma environment in the chapter source.\par\medskip
\noindent\textit{Editorial note.} Pre-open gate: python operations/agent-authoring/tools/precard\_gate.py --source resolve.tex --label lemma-dominate-by-normalized-blowing-up --start 1276 --end 1387 --local-id 28072 returned RESULT=OK (L28072 free, occupancy FREE, public interval clear). Pre-assembly tag self-check: canonical tag 0BBT, not present in reference/public-items.jsonl.\par\medskip
\noindent\textit{Editorial note.} The statement and proof are copied from the pinned chapter file; the cited inputs are kept as external pointers with their canonical Stacks tags.\par\medskip
\noindent\textit{Allowed standard prerequisites (editorial).}\begin{itemize}
\item Isomorphism locus of a modification of normal surfaces. Varieties, the lemma that a birational proper morphism between normal surfaces is an isomorphism over every point of codimension at most one and over every point with finite fibre. Supplies that only finitely many closed points of the target are bad, and the transcendence degree computation for the residue field of a bad curve.
\item Geometric connectedness of the fibres of a modification of normal surfaces. More on Morphisms, the lemma that the fibres of a proper birational morphism towards a normal surface are geometrically connected of dimension at most one. Supplies the description of the fibres over the bad points, hence the finiteness of the set of bad curves.
\item Properties of normalization and the dimension formula for surfaces. Morphisms, the sections on Nagata schemes, the dimension formula and normalization, used without further mention in the proof. Supply the normalization steps and the dimension bookkeeping for surfaces.
\end{itemize}\medskip
\ifdefined\hypertarget\hypertarget{sourcepointer-1}{}\fi
\noindent\textit{Source reference.} Original lemma 0AY8: sources/stacks-chow/tags.txt, source lines 13781--13781. Prerequisite pointer invoked at proof line(s) 1315; it names a standard fact of a sibling chapter that the counted argument uses, and every step of the counted proof built on it is retained verbatim in the delivered body.\par\medskip
\ifdefined\hypertarget\hypertarget{sourcepointer-2}{}\fi
\noindent\textit{Source reference.} Original section 02JT: sources/stacks-chow/tags.txt, source lines 3348--3348. Prerequisite pointer invoked at proof line(s) 1299; it names a standard fact of a sibling chapter that the counted argument uses, and every step of the counted proof built on it is retained verbatim in the delivered body.\par\medskip
\ifdefined\hypertarget\hypertarget{sourcepointer-3}{}\fi
\noindent\textit{Source reference.} Original section 0359: sources/stacks-chow/tags.txt, source lines 4120--4120. Prerequisite pointer invoked at proof line(s) 1298; it names a standard fact of a sibling chapter that the counted argument uses, and every step of the counted proof built on it is retained verbatim in the delivered body.\par\medskip
\ifdefined\hypertarget\hypertarget{sourcepointer-4}{}\fi
\noindent\textit{Source reference.} Original section 035E: sources/stacks-chow/tags.txt, source lines 4125--4125. Prerequisite pointer invoked at proof line(s) 1300; it names a standard fact of a sibling chapter that the counted argument uses, and every step of the counted proof built on it is retained verbatim in the delivered body.\par\medskip
\ifdefined\hypertarget\hypertarget{sourcepointer-5}{}\fi
\noindent\textit{Source reference.} Original lemma 0A21: sources/stacks-chow/tags.txt, source lines 12689--12689. Prerequisite pointer invoked at proof line(s) 1356; it names a standard fact of a sibling chapter that the counted argument uses, and every step of the counted proof built on it is retained verbatim in the delivered body.\par\medskip
\ifdefined\hypertarget\hypertarget{sourcepointer-6}{}\fi
\noindent\textit{Source reference.} Original lemma 0BFP: sources/stacks-chow/tags.txt, source lines 14392--14392. Prerequisite pointer invoked at proof line(s) 1309; it names a standard fact of a sibling chapter that the counted argument uses, and every step of the counted proof built on it is retained verbatim in the delivered body.\par\medskip
\begin{lemma}
\label{lemma-dominate-by-normalized-blowing-up}
Let $X$ be a scheme which is Noetherian, Nagata, and has dimension $2$.
Let $f : Y \to X$ be a proper birational morphism.
Then there exists a commutative diagram
$$
\xymatrix{
X_n \ar[r] \ar[d] &
X_{n - 1} \ar[r] &
\ldots \ar[r] &
X_1 \ar[r] &
X_0 \ar[d] \\
Y \ar[rrrr]  & & & & X
}
$$
where $X_0 \to X$ is the normalization and
where $X_{i + 1} \to X_i$ is the normalized blowing up of $X_i$ at a closed
point.
\end{lemma}

\begin{proof}
We will use the results of Morphisms, Sections
\ref{morphisms-section-nagata},
\ref{morphisms-section-dimension-formula}, and
\ref{morphisms-section-normalization} without further mention.
We may replace $Y$ by its normalization. Let $X_0 \to X$
be the normalization. The morphism $Y \to X$ factors through $X_0$.
Thus we may assume that both $X$ and $Y$ are normal.

\medskip\noindent
Assume $X$ and $Y$ are normal. The morphism $f : Y \to X$ is an isomorphism
over an open which contains every point of codimension $0$ and $1$ in $Y$ and
every point of $Y$ over which the fibre is finite, see Varieties, Lemma
\ref{varieties-lemma-modification-normal-iso-over-codimension-1}.
Hence there is a finite set of closed points $T \subset X$
such that $f$ is an isomorphism over $X \setminus T$. For each $x \in T$
the fibre $Y_x$ is a proper geometrically connected scheme of dimension $1$
over $\kappa(x)$, see
More on Morphisms, Lemma
\ref{more-morphisms-lemma-geometrically-connected-fibres-towards-normal}.
Thus
$$
BadCurves(f) = \{C \subset Y\text{ closed} \mid
\dim(C) = 1, f(C) = \text{a point}\}
$$
is a finite set. We will prove the lemma by induction on the number
of elements of $BadCurves(f)$. The base case is the case where $BadCurves(f)$
is empty, and in that case $f$ is an isomorphism.

\medskip\noindent
Fix $x \in T$. Let $X' \to X$ be the normalized blowup of $X$ at $x$ and let
$Y'$ be the normalization of $Y \times_X X'$. Picture
$$
\xymatrix{
Y' \ar[r]_{f'} \ar[d] & X' \ar[d] \\
Y \ar[r]^f & X
}
$$
Let $x' \in X'$ be a closed point lying over $x$ such that
the fibre $Y'_{x'}$ has dimension $\geq 1$. Let $C' \subset Y'$
be an irreducible component of $Y'_{x'}$, i.e., $C' \in BadCurves(f')$.
Since $Y' \to Y \times_X X'$ is finite we see that $C'$ must map
to an irreducible component $C \subset Y_x$.
If is clear that $C \in BadCurves(f)$.
Since $Y' \to Y$ is birational and hence an isomorphism over points of
codimension $1$ in $Y$, we see that we obtain an injective map
$$
BadCurves(f') \longrightarrow BadCurves(f)
$$
Thus it suffices to show that after a finite number of these
normalized blowups we get rid at of at least one of the bad
curves, i.e., the displayed map is not surjective.

\medskip\noindent
We will get rid of a bad curve using an argument due to Zariski.
Pick $C \in BadCurves(f)$ lying over our $x$. Denote $\mathcal{O}_{Y, C}$
the local ring of $Y$ at the generic point of $C$. Choose an element
$u \in \mathcal{O}_{X, C}$ whose image in the residue field
$R(C)$ is transcendental over $\kappa(x)$ (we can do this because
$R(C)$ has transcendence degree $1$ over $\kappa(x)$ by
Varieties, Lemma \ref{varieties-lemma-dimension-locally-algebraic}).
We can write $u = a/b$ with $a, b \in \mathcal{O}_{X, x}$ as
$\mathcal{O}_{Y, C}$ and $\mathcal{O}_{X, x}$ have the same
fraction fields. By our choice of $u$ it must be the case that
$a, b \in \mathfrak m_x$. Hence
$$
N_{u, a, b} = \min
\{\text{ord}_{\mathcal{O}_{Y, C}}(a), \text{ord}_{\mathcal{O}_{Y, C}}(b)\} > 0
$$
Thus we can do descending induction on this integer.
Let $X' \to X$ be the normalized blowing up of $x$
and let $Y'$ be the normalization of $X' \times_X Y$ as above. We will
show that if $C$ is the image of some bad curve $C' \subset Y'$
lying over $x' \in X'$, then
there exists a choice of $a', b' \mathcal{O}_{X', x'}$
such that $N_{u, a', b'} < N_{u, a, b}$. This will finish the proof.
Namely, since $X' \to X$ factors through the blowing up, we see that
there exists a nonzero element $d \in \mathfrak m_{x'}$ such that
$a = a' d$ and $b = b' d$ (namely, take $d$ to be the local equation
for the exceptional divisor of the blowup). Since $Y' \to Y$
is an isomorphism over an open containing the generic point of $C$
(seen above) we see that $\mathcal{O}_{Y', C'} = \mathcal{O}_{Y, C}$.
Hence
$$
\text{ord}_{\mathcal{O}_{Y, C}}(a) =
\text{ord}_{\mathcal{O}_{Y', C'}}(a' d) =
\text{ord}_{\mathcal{O}_{Y', C'}}(a') +
\text{ord}_{\mathcal{O}_{Y', C'}}(d) >
\text{ord}_{\mathcal{O}_{Y', C'}}(a')
$$
Similarly for $b$ and the proof is complete.
\end{proof}

\end{document}
